\chapter{Introduction}\label{chapter:intro}

Causal inference methods in the social and biomedical sciences are often motivated by the promise 
that these methods can identify and isolate the ``effects of causes''~\citep{holland1986}. 
However, we are often interested not in the total effect of an intervention but instead in 
the mechanisms through which an intervention induces an effect. This is because there are 
many situations in which an intervention --- such as a social policy or a drug --- affects 
an outcome in multiple, sometimes conflicting ways. A canonical example comes from~\cite{hesslow1976}, 
which discusses a randomized experiment that tested whether a new birth control pill increased the 
risk of blood clots in women. As discussed in~\cite{pearl2001}, this drug could have a direct effect 
on the rate of blood clots in women by causing clots to form in some women who would otherwise not 
have been at risk. However, the drug could also operate through a second pathway whose sign opposes 
the direct effect: the birth control pill could indirectly reduce the risk of developing blood clots 
by reducing pregnancies, which are known to lead to higher rates of blood clots in women. A simple 
causal inference framing asks only: does taking the drug increase the risk of blood clots for women? 
However, this question does not address our concerns with the drug's safety: it is entirely possible 
that this drug causes blood clots in some women, but because the women who take the drug get pregnant 
at a lower rate, there is no observable difference in blood clot rates for women who did and did not 
take the drug. The policy-relevant question is what the direct effect of the drug on blood clot risk 
is, after accounting for the indirect effect that operates through the lowered risk of pregnancy. 

Mediation analysis attempts to elucidate the mechanisms through which an intervention affects an 
outcome of interest. It does so by decomposing the total causal effect of an intervention on an 
outcome into the causal effects of distinct pathways through which the intervention affects the 
outcome. Informally, the total observable causal effect of an intervention on an outcome is typically 
decomposed into a \textbf{direct effect} that represents the unmediated effect that an intervention 
has on the outcome and an \textbf{indirect effect} that represents the effect that an intervention 
has on an outcome only through the way that it affects mediator variables that themselves affect 
the outcome. 

Since statisticians first became interested in formalizing mediation analysis methods in the 1990s 
and early 2000s, several different versions of mediational effects have been proposed in the literature,
each with their attendant interpretations, required assumptions, and theoretical issues. Natural effects, 
which capture the effects of an intervention through a mediator compared to the 
effect of the intervention had the mediator been set at the level it would have ``naturally'' 
attained under no intervention, were the first class of effects proposed in the literature
~\citep{robinsgreenland1992, pearl2001, avinetal2005}. While theoretically appealing, these 
effects have been controversial because they require an empirically untestable
assumption for identification, and the practical usefulness of their implications has also 
been called into question~\citep{vanderweelevansteelandt2009}. 
A recently developed alternative, interventional 
effects, require significantly weaker assumptions. Interventional effects estimate the 
effect of joint interventions on the exposure and the mediator that set the level of the mediator to a 
random draw from the distribution of the mediator from a group that either did or did not receive 
treatment. They have the advantage of not requiring the 
untestable assumptions~\citep{vanderweeleetal2014}. 

With the popularization of interventional effects, recent research in mediation analysis has also 
attempted to extend the interventional effects framework to problems with multiple 
mediators~\citep{vanderweeletchetgentchetgen2017, linvanderweele2017, vansteelandtdaniel2017}. This has 
been seen as an important task because most mediation analysis of practical interest contain more than 
one mediator~\citep{vansteelandtdaniel2017}. Additionally, natural effect definitions face many issues 
when dealing with multiple mediators in most situations: even under strong assumptions, it is not 
possible to identify the causal pathways for natural effects when there exists dependence between
the mediators~\citep{avinetal2005, vanderweeleetal2014}. 

Despite the appeal of interventional effects definitions, recent work has argued that such effects do 
not carry the same mediational interpretation that natural effects do~\citep{miles2023}. In this 
dissertation, I will further develop this insight by analyzing natural and interventional 
effects in terms of a conceptual distinction between descriptive and prescriptive interpretations of 
effects~\citep{pearl2001}.~\cite{pearl2001} claims that descriptive effects concerns the action of 
variables, such as mediators, under ``natural'' conditions where such variables are not being directly
intervened upon. These descriptive effects can have policy implications for exactly the reason why 
we would care about the direct and indirect mechanisms in the birth control example. Prescriptive 
effects, meanwhile, give causal effects under exact hypothetical or actual experimental conditions
and concern the causal effect of a particular intervention rather than the descriptive properties 
of causal pathways in the causal system. For example, a joint intervention where we assign both 
the birth control pill and a blood thinner to a patient may be tried, and we would be interested in 
if this combination reduces blood clots regardless of the causal pathways at work. 

Given that interventional effects are easier to identify, interventional effects have recently 
become particularly popular in applied work as conceptual substitutes for natural 
effects~\citep{miles2023}. Additionally, as I will show in this dissertation, the descriptive 
intent behind natural effects, which motivates an emphasis on features such as decomposing the total 
effect into path-specific effects, have also seemed to influence the development of recent
interventional effect definitions for multiple mediator problems. The results
of~\cite{miles2023} reinforce this division and argue that interventional effects cannot be 
interpreted in the way natural effects are. Natural effect estimation is motivated by the 
descriptive goal of describing mediating pathways, while interventional effects are generally 
motivated by the prescriptive goal of estimating the causal 
effects of particular stochastic interventions on the distribution of the mediator. Therefore,
interventional effects must be interpreted prescriptively. 

In multiple mediator problems, natural effects are plagued by impassable barriers to identification. 
I will argue that interventional effects, while requiring looser identification assumptions, 
are also unappealing targets in such problems because of the difficulty of finding practical 
situations where such effects would be useful quantities to estimate. This stands in contrast to 
targets like controlled effects and even interventional effects in one mediator problems, where the 
practical relevance of the interventions these effects map onto is more tractable. 

If neither natural nor interventional effects are appealing targets in multiple mediator problems, 
what types of effects can be productively pursued for such problems? Following the suggestion
of~\cite{miles2023}, I argue that a version of interventional effects first proposed
by~\cite{vanderweelerobinson2014} holds important practical importance. These effects are developed
for problems where the exposure variable is considered to be non-manipulable 
(e.g. race, sex, social class) and aim to estimate the effect of hypothetical interventions to 
equalize the distribution of a mediator between two populations with different values of the 
``exposure''. Such effects have been productively developed in the health disparities 
literature~\citep{jacksonvanderweele2018,jackson2021}. I relate the formal framework I use for 
natural and interventional effects to these effects and use the results of~\citep{jacksonvanderweele2018} 
for single mediator problems with confounding to propose extensions of these effects to two mediator 
problems. 

Finally, I demonstrate the practical usefulness and interpretation of these effects 
through an illustrative application in which I use data from~\cite{kuhaetal2021} to estimate 
these effects for a problem in which intergenerational social class mobility is mediated by intelligence
and education. To do this, I develop code to implement a general Monte Carlo estimation strategy that can be
deployed to flexibly estimate different versions of interventional effects for multiple mediators. The code 
notably separates the estimation of the conditional distributions for the mediators and outcome --- I use
linear models in this paper, but this does not have to be so --- from the Monte Carlo estimation procedure 
that computes interventional effects. The idea of using Monte Carlo estimation for mediation analysis has 
precursors~\citep{imaietal2010b,vansteelandtdaniel2017}, but previous methods have not been developed for 
VanderWeele-Robinson effects, and I am also not aware of whether they are able to be used to estimate other 
interventional effect types. 

Section 2 of this paper gives a brief history of the mediation analysis literature and overview of other 
active areas of research. Section 3 introduces the potential outcomes framework for causal 
inference~\citep{rubin1974}, which serves as the formal framework for the rest of the dissertation. 
Section 4 develops the theory, assumptions, and interpretation of natural effects for one and two mediator
problems. Section 5 develops the theory, assumptions, and interpretation of interventional effects for one 
and two mediator problems and additionally develops the theory for VanderWeele-Robinson interventional 
effects with one and two mediators. Section 6 applies the VanderWeele-Robinson interventional 
effects developed in Section 5 to analyze the aforementioned social class mobility dataset. Section 7
concludes. 